% questions.tex —— 1985 年·上海卷·文史类 · 题目内容片段（无答案）
% 由 tools/generate_projects.py 依来源总账自 1985/1985_数学真题_上海文_试题.md 生成，勿手改（改源 MD 后重跑）


\begin{gktip}
\textbf{来源说明}：本题面依据固定源 PDF《1985 年上海市普通高等学校招生统一考试题目·数学（文史类）》第 1—2 页逐题转录。原卷页码为 4-37、4-38；该 PDF 不含答案页，下列答案和解析由本库独立推导，不声称代表官方发布版本。

\textbf{图形说明}：第 3 题明确要求画出两条曲线，LaTeX 版依据方程重绘为 TikZ 矢量图。第 5、8 题题干写有“如图”，但固定源两页中未见可独立辨认的对应图形；本文保留可见文字，不凭空补画源图。
\end{gktip}

\begin{question}
\gkti{1}（本题满分 30 分）

本题共有 10 个小题，每一个小题满分 3 分，只要求直接填写结果。

\Needspace{5\baselineskip}
（1） 不等式

\nopagebreak[4]
\[
|x+2|<1
\]

的解集是　　　　。

\Needspace{5\baselineskip}
（2） 函数

\nopagebreak[4]
\[
y=\sqrt[8]{x}-1
\]

的反函数是　　　　。

\Needspace{5\baselineskip}
（3） 函数

\nopagebreak[4]
\[
y=\operatorname{tg}2x
\]

的最小正周期是　　　　。

（4） 设等差数列 $\{a_n\}$ 中，$a_1=2$，$a_2=5$，则
$a_{10}=$　　　　。

\Needspace{5\baselineskip}
（5） 方程

\nopagebreak[4]
\[
9^x=3^{x+1}
\]

的解是　　　　。

\Needspace{5\baselineskip}
（6） 用不等号“$>$”或“$<$”连接：

\nopagebreak[4]
\[
\left(\frac12\right)^{\frac17}\ \underline{\hspace{2em}}\ 1.
\]

（7） 如果一个圆的圆心在点 $(2,4)$，并且经过点 $(0,3)$，则这个圆的方程是　　　　。

\Needspace{5\baselineskip}
（8） 直线

\nopagebreak[4]
\[
y=-x+3
\]

的倾斜角的大小是　　　　。

\Needspace{5\baselineskip}
（9） 设 $\lg2=a$，则

\nopagebreak[4]
\[
\log_2 25=
\]

　　　　。

（10） 已知 $O$ 为直角坐标系的原点，$A$ 点的坐标为 $(1,3)$，$B$ 点的坐标为 $(0,3)$，若把 $\triangle OAB$ 绕 $y$ 轴旋转一周，则所得旋转体的体积是　　　　。

\end{question}
\begin{question}
\gkti{2}（本题满分 15 分）

本题共有 5 个小题，每一个小题都给出代号为 A、B、C、D 的四个结论，其中只有一个结论是正确的。把正确结论的代号写在题后的圆括号内，选对得 3 分，不选、选错或者选出的代号超过一个（不论是否都写在圆括号内），一律得 0 分。

\Needspace{5\baselineskip}
（1） 函数

\nopagebreak[4]
\[
y=-\sin2x+3
\]

的最大值是\mbox{（　　）}

（A）2　　（B）3　　（C）4　　（D）5

（2） 设复数 $z_1=1+i$，$z_2=1-3i$，则复数
$z=z_1z_2$ 在复平面内所表示的点位于\mbox{（　　）}

（A）第一象限　　（B）第二象限　　（C）第三象限　　（D）第四象限

\Needspace{5\baselineskip}
（3） 已知函数

\nopagebreak[4]
\[
f(x)=\lg x^2
\]

\Needspace{5\baselineskip}
的定义域为 $F$，函数

\nopagebreak[4]
\[
g(x)=2\lg x
\]

的定义域为 $G$，那么\mbox{（　　）}

（A）$F\cap G=\varnothing$　　（B）$F=G$

（C）$F\subset G$　　（D）$G\subset F$

（4） 若一个四棱锥的底面是凸四边形，它的顶点到底面各边的距离都相等，则这个棱锥的底面四边形\mbox{（　　）}

（A）必为正方形　　（B）必有内切圆

（C）必有外接圆　　（D）必既有内切圆又有外接圆

（5） $a>0$ 且 $b>0$，是 $ab>0$ 的\mbox{（　　）}

（A）充分但不必要的条件　　（B）必要但不充分的条件

（C）充分而且必要的条件　　（D）既不充分又不必要的条件

\Needspace{0.4\textheight}
\end{question}
\begin{question}
\gkti{3}（本题满分 9 分）

本题共有 2 个小题，第（1）小题满分 4 分，第（2）小题满分 5 分。

在平面直角坐标系内，画出下列方程的曲线（只要求画出图形）：

\Needspace{5\baselineskip}
（1）　

\nopagebreak[4]
\[
\frac{x}{3}-\frac{y}{2}=-1;
\]

\Needspace{5\baselineskip}
（2）　

\nopagebreak[4]
\[
\frac{x^2}{9}-\frac{y^2}{4}=-1.
\]
\nopagebreak[4]
\begin{center}% 1985 年上海卷文史类第 3 题：直线与竖轴双曲线
\begin{tikzpicture}[scale=0.72,>=stealth]
  \begin{scope}[shift={(-3.0,0)}]
    \draw[->] (-2.1,0)--(2.25,0) node[right] {$x$};
    \draw[->] (0,-1.0)--(0,2.8) node[above] {$y$};
    \draw[thick] (-2.0,0.67)--(1.2,2.8);
    \fill (0,2) circle (1.1pt);
    \node[below left] at (0,0) {$O$};
    \node[below] at (0,-1.25) {（1）};
  \end{scope}
  \begin{scope}[shift={(2.0,0)}]
    \draw[->] (-1.85,0)--(1.9,0) node[right] {$x$};
    \draw[->] (0,-2.4)--(0,2.4) node[above] {$y$};
    \draw[thick,domain=-1.65:1.65,samples=100]
      plot (\x,{2*sqrt(1+\x*\x/9)});
    \draw[thick,domain=-1.65:1.65,samples=100]
      plot (\x,{-2*sqrt(1+\x*\x/9)});
    \fill (0,2) circle (1.1pt);
    \fill (0,-2) circle (1.1pt);
    \node[below left] at (0,0) {$O$};
    \node[below] at (0,-2.75) {（2）};
  \end{scope}
\end{tikzpicture}
\end{center}
\nopagebreak[4]
\end{question}
\begin{question}
\gkti{4}（本题满分 20 分）

本题共有 3 个小题，第（1）小题满分 5 分，第（2）小题满分 7 分，第（3）小题满分 8 分。

\Needspace{5\baselineskip}
（1） 求

\nopagebreak[4]
\[
(2+x^3)^8
\]

的展开式中 $x^{10}$ 的系数。

（2） 从六个数字 $1,2,3,5,7,9$ 中任取四个不同的数字，有多少种取法？由这六个数字可以组成多少个没有重复数字的四位偶数？

（3） 在一列球中，第 1 个球的半径为 1，第 2 个球的直径是第 1 个球的半径，以后每一个球的直径都是前一个球的半径，这样无限继续下去，求所有这些球的表面积的和。

\end{question}
\begin{question}
\gkti{5}（本题满分 10 分）

如图，直线 $a\parallel$ 直线 $b$，点 $P$ 在 $a\text{、}b$ 所确定的平面外，$PA\perp a$，且交 $a$ 于 $A$ 点；$AB\perp b$，且交 $b$ 于 $B$ 点。

求证：$PB\perp b$。

\begin{gktip}
固定源页面保留“如图”题干，但未见可独立辨认的配图；证明仅依据可见直线关系转录。
\end{gktip}

\end{question}
\begin{question}
\gkti{6}（本题满分 12 分）

本题共有 2 个小题，第（1）小题满分 4 分，第（2）小题满分 8 分。

\Needspace{5\baselineskip}
（1） 证明万能公式：

\nopagebreak[4]
\[
\sin\alpha=\frac{2\operatorname{tg}\frac\alpha2}
{1+\operatorname{tg}^2\frac\alpha2}.
\]

\Needspace{5\baselineskip}
（2） 已知

\nopagebreak[4]
\[
\cos\alpha=-\frac13,\qquad
\pi<\alpha<\frac{3\pi}{2},
\]

\Needspace{5\baselineskip}
求

\nopagebreak[4]
\[
\cos\left(\frac\alpha2+\frac{5\pi}{6}\right).
\]

\end{question}
\begin{question}
\gkti{7}（本题满分 10 分）

设 $\triangle ABC$ 的两个内角 $A\text{、}B$ 所对的边分别为 $a\text{、}b$。若复数 $z_1$ 的模为 $a$，幅角为 $B$；复数 $z_2$ 的模为 $b$，幅角为 $-A$。求证：$z_1+z_2$ 为实数。

\end{question}
\begin{question}
\gkti{8}（本题满分 14 分）

\Needspace{5\baselineskip}
如图，由抛物线

\nopagebreak[4]
\[
y=x^2+2
\]

\Needspace{5\baselineskip}
上的点 $M(x_0,y_0)$ 向直线

\nopagebreak[4]
\[
y=\frac12x
\]

\Needspace{5\baselineskip}
作垂线，垂足为 $N$，延长 $MN$ 至 $P$ 点，使点 $N$ 分线段 $MP$ 成定比

\nopagebreak[4]
\[
\lambda=\frac{MN}{NP}=4.
\]

（1） 用 $x_0,y_0$ 表示点 $N$ 的坐标 $(x_1,y_1)$；

（2） 用 $x_0,y_0$ 和 $x_1,y_1$ 表示点 $P$ 的坐标 $(x,y)$；

（3） 求：当点 $M$ 沿抛物线移动时，点 $P$ 的轨迹方程，并说明轨迹是什么图形。

\begin{gktip}
固定源页面保留“如图”题干，但未见可独立辨认的配图；本题坐标推导仅依据可见文字转录。
\end{gktip}

\end{question}